\documentclass[11pt]{article}
\usepackage[margin=1.15in]{geometry}
\usepackage{amsmath,amssymb,amsthm,mathtools}
\usepackage{xcolor}
\usepackage[colorlinks=true,linkcolor=blue!60!black,citecolor=blue!60!black,urlcolor=blue!60!black]{hyperref}

\newtheorem{theorem}{Theorem}[section]
\newtheorem{proposition}[theorem]{Proposition}
\newtheorem{lemma}[theorem]{Lemma}
\newtheorem{corollary}[theorem]{Corollary}
\theoremstyle{definition}
\newtheorem{definition}[theorem]{Definition}
\theoremstyle{remark}
\newtheorem{remark}[theorem]{Remark}

\newcommand{\C}{\mathbb{C}}
\newcommand{\R}{\mathbb{R}}
\newcommand{\Z}{\mathbb{Z}}
\newcommand{\w}{\omega}
\newcommand{\Sp}{\mathrm{Sp}}
\newcommand{\dg}{\dagger}
\newcommand{\WF}{\operatorname{WF}}
\newcommand{\tr}{\operatorname{tr}}

\title{The Universal Vacuum of the Fourth-Order Scalar Field:\\
Metric Orbits, Fock Sectors, and the Krein Boundary}
\author{GPT-5.6.sol\thanks{OpenAI model and principal programme author.}
\and Asger Alstrup Palm\thanks{Programme orchestrator and corresponding human contact; Honorary Professor, DTU Compute, Technical University of Denmark. \texttt{asger@area9.dk}.}}
\date{Open-repository preprint, July 2026}

\begin{document}
\maketitle

\begin{abstract}
For the free fourth-order scalar field
$(\Box+m_1^2)(\Box+m_2^2)\phi=0$ we separate three structures that are
commonly conflated: the geometry of positive metric operators, the geometry
of vacuum functionals, and the geometry of the dynamics.  Modewise, the
Bender--Mannheim generator is a complexified beam-splitter generator in
canonical normal-form coordinates --- it carries divergent Cartan
distortion while creating no quanta there --- and no symplectic frame
tames its singular values: particle content is a vacuum-ray functional,
not a Cartan-norm functional.  The entire family of admissible positive
metrics fixes the physical vacuum ray and induces one and the same
selected complex covariance --- the \emph{universal selected quasifree
functional} within the stated real-form class; it carries exactly $1/3$
of a quantum per ultraviolet mode pair relative to the two-field Fock
vacuum, and, on the common auxiliary canonical algebra, its
representation is disjoint from the two-field Fock representation in
every spatial dimension --- an obstruction no choice of metric can cure.
The ultraviolet Fock-sector hierarchy of selected vacua terminates: no
ultraviolet invariant for $d\le3$, the invariant $m_1^2+m_2^2$ for $4\le
d<8$, the unordered mass pair for $d\ge8$; global equivalence involving
the doubly massless vacuum additionally requires an infrared analysis,
which we identify but do not perform.  The selected vacuum's Wightman
function is exactly the positive-energy spectral splitting of the
fourth-order commutator for every mass splitting; its confluent limit is
the massive analogue of the Bateman--Turok Krein covariance, agreeing
with their doubly massless theory at $m=0$ up to an action-sign
reversal: metric selection and the spectral condition determine the same
complex quasifree covariance, the quantizations differing in involution
and completion --- positive pseudo-Hermitian for split masses,
Krein/Jordan at resonance.  Finally the selected functional satisfies a
fourth-order microlocal spectrum condition: its wavefront set is the
standard positive-frequency Hadamard set, its short-distance singularity
is $V\log\sigma_+$ with smooth $V$ of universal diagonal value
$1/(16\pi^2)$ (remainder $O(\rho^2\log\rho)$, not smooth), and its
restriction to the local partial-fraction fields is
$(\pm)$Klein--Gordon--Hadamard.  All identities are machine-verified.
\end{abstract}

\section*{AI authorship and computational accountability}
GPT-5.6.sol developed the derivations, computational checks, claim audit, and
manuscript.  Asger Alstrup Palm commissioned and orchestrated the programme
and is the corresponding human contact.  The project is an experiment in
whether a non-specialist human, using AI systems as research collaborators,
can organize falsifiable computer-assisted mathematical physics.  The exact
status of each repository check is stated where it is used, and no check is
promoted beyond its documented scope.  The manuscript remains subject to
expert review.

\section{Introduction}\label{sec:intro}

The Pais--Uhlenbeck (PU) oscillator and its field-theoretic parent, the
free fourth-order scalar
\begin{equation}\label{eq:eom}
(\Box+m_1^2)(\Box+m_2^2)\,\phi=0,\qquad m_1>m_2\ge0,
\end{equation}
admit two apparently different resolutions of the Ostrogradsky ghost
problem.  In the PT-symmetric approach of Bender and Mannheim
\cite{BM2008PRL,BM2008PRD}, each mode carries a positive metric operator
$\eta=e^{-Q}$ that renders the dynamics unitary on a Hilbert space; in the
Krein-space approach, recently developed for the pure fourth-order theory
$\Box^2\phi=0$ by Bateman and Turok \cite{BT2026}, the state space is
indefinite and probabilities are controlled by a hidden ghost-parity
symmetry.  Two companion papers \cite{companion1,companion2} analyzed the
modewise PT construction as finite-dimensional symplectic geometry: the
positive metric is unique up to a stabilizer orbit, is selected
variationally by minimum Cartan distortion, and the modewise theory is
exactly solvable in all the respects used below.  The minimum-distortion
selection is itself a Cartan-projection theorem, $F(S)=2\|\mu(S)\|^2$
with $\mu(S_+)=(r,r)$ on the $C_2$ Weyl-chamber wall --- but not a
routine nearest-point (Mostow/Kempf--Ness) application: the stabilizer
$SO(2,\C)^2$ is not compatible with the Cartan involution
($H\cap K\cong\Z_2\times\Z_2$), and the proof passes through the
$\theta$-compatible hull $SL(2,\C)\times SL(2,\C)$ \cite{companion2}.

This paper carries that analysis to the field theory, where three
different geometries are in play and must not be conflated:
\begin{equation}\label{eq:threegeom}
\text{metric geometry}\ \neq\ \text{vacuum geometry}\ \neq\
\text{dynamical geometry}.
\end{equation}
The metric geometry is the symmetric-space geometry of positive operators,
$G/K$ for $G=\Sp(2n,\C)$, where the Cartan projection $\mu$ measures
nonunitarity \cite{companion2}.  The vacuum geometry is the geometry of
Gaussian vacuum rays, which is governed by a different quotient: the
stabilizer $P_\Omega=\{g:\hat g\,\Omega\in\C^\times\Omega\}$ of the
vacuum line is the Siegel parabolic of the vacuum Lagrangian in the
oscillator representation, and vacuum rays live in $G/P_\Omega$.  The
vacuum-ray map $q_\Omega:G\to G/P_\Omega$ does not factor through
$q_K:G\to G/K$ --- no $G$-equivariant $F$ with $q_\Omega=F\circ q_K$
exists, because the compact $K\cong\Sp(n)$ does not preserve the vacuum
ray (only $K\cap P_\Omega\cong U(n)$ does), so $q_\Omega$ is not
constant on $K$-cosets; the correct geometry is the correspondence
$G/K\leftarrow G\rightarrow G/P_\Omega$ of two inequivalent quotient
maps.  The dynamical geometry is the spectral structure of the
Hamiltonian, which degenerates (Jordan blocks) exactly at the resonant
point $m_1=m_2$ \cite{companion1}.  Our results, in brief:

\begin{enumerate}
\item \emph{Cartan--parabolic decoupling}
(Section~\ref{sec:decoupling}).  In canonical normal-form coordinates the
Bender--Mannheim generator is a complexified beam-splitter generator,
$Q'=r\,(a_1a_2^{\dg}+a_1^{\dg}a_2)$: it lies in the complexified
number-preserving (Levi) algebra of $P_\Omega$ --- $e^{-Q'/2}$ with real
$Q'$ is \emph{not} a unitary $SU(2)$ beam splitter --- fixes the
canonical vacuum exactly, and yet has Cartan projection $(r,r)$ with
$r=\log\frac{\w_1+\w_2}{\w_1-\w_2}\to\infty$ in the ultraviolet.
Moreover no symplectic frame tames the singular values.  Hence particle
content is not a function of the Cartan projection, and no coercive
lower bound $\text{vacuum excitation}\ge f(\|\mu\|)$ with
$f(r)\to\infty$ can hold: large Cartan distortion does not force
particle content.  Geometrically, the ultraviolet limit is escape to
infinity along the $C_2$ chamber wall \emph{within a single vacuum-ray
fibre} of $q_\Omega$.

\item \emph{Orbit constancy and a universal obstruction}
(Section~\ref{sec:obstruction}).  The physical annihilator covectors are
eigenvectors of the stabilizer generators, so the entire admissible
metric family --- the Gaussian orbit and all $W=f(N_1,N_2)>0$ ---
fixes the physical vacuum ray.  The vacuum functional is therefore
metric-independent.  Its occupation relative to the two-field Fock vacuum
is exactly
\begin{equation}\label{eq:N}
\langle N\rangle(k)
=\frac{\w_2(\w_1^2+\w_2^2)}{2\w_1(\w_1^2+\w_1\w_2+\w_2^2)}
\ \xrightarrow[k\to\infty]{}\ \frac13 ,
\end{equation}
a universal constant; the total relative quanta scale as
$\Theta(V\Lambda^d)$ and the representations are disjoint in every
spatial dimension.  No choice of metric can cure this.

\item \emph{Terminating ultraviolet sector hierarchy}
(Section~\ref{sec:sectors}).  Selected vacua of different mass pairs
satisfy, in common field variables,
\begin{equation}\label{eq:sector}
1-f(k)=\frac{(\Sigma-\bar\Sigma)^2}{12\,k^4}+O(k^{-6}),
\qquad \Sigma=m_1^2+m_2^2 .
\end{equation}
At matched $\Sigma$ the $k^{-6}$ coefficient also vanishes, and the
first remaining obstruction is
$\tfrac{29}{576}\,(\delta^2-\bar\delta^2)^2k^{-8}$,
$\delta=\tfrac12(m_1^2-m_2^2)$; the invariants determine the mass pair
only up to interchange, $(m_1^2,m_2^2)\sim(m_2^2,m_1^2)$.  Hence the
\emph{ultraviolet} hierarchy: no UV invariant for $d\le3$; $\Sigma$
for $4\le d<8$; the unordered pair $(\Sigma,\Pi)$ for $d\ge8$; and the
classification terminates.  Global (GNS) equivalence involving the
doubly massless vacuum requires a separate infrared analysis: against
the $\Box^2$ anchor the relative occupation diverges as
$(m_1{+}m_2)m_1m_2/(6k^3)$ in the infrared, and
$\int_0k^{d-1}N_{\mathrm{rel}}\,dk$ diverges precisely for $d\le3$.

\item \emph{The spectral bridge} (Section~\ref{sec:bridge}).  In the
action convention of \eqref{eq:action}, the selected vacuum's Wightman
function is exactly the positive-energy spectral splitting of the
fourth-order commutator,
\begin{equation}\label{eq:Wspectral}
\widetilde W^+_{12}(p)
=\varepsilon\,\theta(p^0)\,
\frac{\delta(p^2-m_1^2)-\delta(p^2-m_2^2)}{m_1^2-m_2^2},
\qquad
\widetilde W^+_{mm}(p)=-\varepsilon\,\theta(p^0)\,\delta'(p^2-m^2),
\end{equation}
for all $\Delta=m_1^2-m_2^2\ge0$.  The confluent case is a
\emph{massive Bateman--Turok-type} Krein/Jordan covariance; at the
doubly massless point $m=0$, and after matching the action sign and
infrared extension, it agrees with the covariance of \cite{BT2026} ---
in our $\varepsilon=+1$ convention it is its action-sign reversal.
Metric selection and the spectral condition determine the same complex
quasifree covariance; the quantizations differ in involution and
completion --- positive $\eta$-inner product where it exists
($\Delta>0$), Krein space with ghost parity where it must ($\Delta=0$).
The covariance is analytic across the resonant boundary while the metric
and diagonalizability fail there.

\item \emph{Fourth-order microlocal spectrum condition}
(Section~\ref{sec:hadamard}).
$W^+_{12}$ is a bisolution with the correct fourth-order commutator, its
wavefront set is the standard positive-frequency Hadamard set
$\mathcal C^+$, and its short-distance singularity is logarithmic with
universal diagonal coefficient: the $1/\rho^2$ singularity of the
Klein--Gordon factors is mass-independent and cancels, leaving the
Hadamard form $W^+_{12}=V_{12}\log\sigma_++H_{12}$ with $V_{12},H_{12}$
smooth and $V_{12}$ of universal coincidence value; in flat coordinates
\begin{equation}
W^+_{12}(x)=\frac{1}{8\pi^2}\log\rho
+\frac{\Sigma}{64\pi^2}\,\rho^2\log\rho+\dots+C^\infty,
\end{equation}
i.e.\ leading coefficient $1/(16\pi^2)$ for $\log\rho^2$ with a
remainder that is $O(\rho^2\log\rho)$ --- milder than the leading
logarithm but \emph{not} smooth for nonzero masses.  Its restrictions to
the local partial-fraction fields are exactly $(\pm)$Klein--Gordon
Hadamard functions.  The fourth-order parametrix
$H_P=(H_{m_1}-H_{m_2})/(m_1^2-m_2^2)$ provides the local microlocal
input for developing Wick powers and perturbative time-ordered products.
\end{enumerate}

Together these give the corrected physical picture: \emph{the universal
selected vacuum is a fourth-order Hadamard quasifree functional; its
apparently infinite particle content relative to the split two-field
vacuum reflects inequivalent complex structures, not pathological local
ultraviolet singularities} --- and the Krein formulation is not the
remnant of a divergent positive theory but a different completion of the
same, analytically continued covariance.

\paragraph{What ``universal'' means.}
Throughout, the universal object is the \emph{selected complex quasifree
covariance}: it is constant over the entire admissible metric orbit
(Theorem~\ref{thm:constancy}) and analytic in the masses across the
resonant boundary.  It is \emph{not} claimed to be positive under one
common involution across that boundary; its global Fock realization
depends on dimension and infrared behaviour
(Section~\ref{sec:sectors}); ``particle number'' is defined only
relative to a chosen reference complex structure; and the massive
confluent theory is the massive analogue of --- not identical to --- the
doubly massless Bateman--Turok theory.

\paragraph{Verification.}
Every identity below --- the beam-splitter identity, the occupation
\eqref{eq:N}, the eigenvector/constancy lemma, the sector expansions with
their exact coefficients, the Wightman functions and their confluent
limits, the commutator normalizations, the logarithmic coefficient and
the $\rho^2\log\rho$ remainder coefficient, and the infrared divergence
of the relative occupation against the massless anchor ---
is machine-verified symbolically, with independent numerical checks.
The artifact is the root of the \texttt{area9innovation/weyl-gravity}
repository accompanying \cite{companion1,companion2} (scripts
\texttt{symbolic/verify\_paper3\_audit.py} (P1--P10),
\texttt{symbolic/verify\_hadamard.py} (H1--H6),
\texttt{symbolic/verify\_paper3\_referee.py} (R1--R5); SymPy 1.14.0,
Python 3.14; cite the repository commit hash to fix the artifact).

\paragraph{Conventions.}
The theory is defined by the action
\begin{equation}\label{eq:action}
S[\phi]=\frac{\varepsilon}{2}\int
\phi\,(\Box+m_1^2)(\Box+m_2^2)\,\phi,
\qquad \varepsilon=+1\ \text{fixed throughout},
\end{equation}
whose canonical structure fixes the fourth-order commutator
normalization $E_P'''(0)=+\varepsilon$; the equation of motion
\eqref{eq:eom} alone does not fix the symplectic sign, and every signed
object below (commutator, Wightman function, logarithmic coefficient,
partial-fraction signature, comparison with \cite{BT2026}) carries
$\varepsilon$.  We further distinguish three layers: the \emph{complex
covariance} (a bilinear functional on the complex field algebra
generated by the smeared field with no conjugation selected), the
\emph{involution} (a real form of that algebra), and the
\emph{positive/Krein completion} built on the involution; ``same
functional, different completion'' always refers to this hierarchy.
Modewise notation follows \cite{companion1,companion2}: for spatial
momentum $k$, the mode is a PU pair with
$\w_i(k)=\sqrt{k^2+m_i^2}$, canonical variables $V=(x,y,p,q)$ related to
the field variable by $z=iy$ (the PT rotation), normalized coordinates
with $d_xd_y=\w_1\w_2$, rapidity $r(k)=\log\frac{\w_1+\w_2}{\w_1-\w_2}
=\log\frac{(\w_1+\w_2)^2}{\Delta}$, selected diagonalizer
$S_+=e^{(r/2)B}$ and metric $\eta=e^{-Q}$, $Q=\alpha pq+\beta xy$.  We
write $\Delta=m_1^2-m_2^2$, $\Sigma=m_1^2+m_2^2$, $\Pi=m_1^2m_2^2$.
\emph{Cartan convention}: $\mu(S)$ denotes the Gram--Cartan data
$\mu(S)=\log\operatorname{spec}(S^{\dg}S)$ restricted to the closed
chamber, i.e.\ \emph{twice} the standard singular-value projection
$\mu_{\mathrm{std}}(S)=\log s(S)$; the singular values of
$S_+=e^{(r/2)B}$ are $e^{\pm r/2}$, so $\mu(S_+)=(r,r)$ while
$\mu_{\mathrm{std}}(S_+)=(r/2,r/2)$.  This is the convention
synchronized with the distortion dictionary
$F(S)=\|\log S^{\dg}S\|_F^2=2\|\mu(S)\|^2$ of \cite{companion2}.

\section{Metric geometry versus vacuum geometry}\label{sec:decoupling}

\subsection{The beam-splitter identity}

In the canonically normalized coordinates the mode annihilation operators
of the auxiliary canonical vacuum are $a'_j$ with
$x'=(a'_1+a'^{\dg}_1)/\sqrt2$, etc.  A one-line computation (different
modes commute, so no ordering terms arise) gives the exact operator
identity
\begin{equation}\label{eq:beamsplitter}
Q'=r\,(x'y'+p'q')=r\,\big(a'_1a'^{\dg}_2+a'^{\dg}_1a'_2\big).
\end{equation}
Thus $Q'$ lies in the complexified number-preserving beam-splitter
algebra --- the Levi algebra of the vacuum-line parabolic $P_\Omega$;
note that $e^{-Q'/2}$ with real $Q'$ is \emph{not} a unitary $SU(2)$
beam splitter, the generator being Hermitian rather than
anti-Hermitian --- with $[Q',N_{\mathrm{tot}}]=0$ and
$Q'\,\Omega_{\mathrm{can}}=0$, so
\begin{equation}
e^{-Q'/2}\,\Omega_{\mathrm{can}}=\Omega_{\mathrm{can}}
\qquad\text{exactly},
\end{equation}
even though the Cartan projection of the implemented symplectic map is
$\mu(S_+)=(r,r)$ with $r(k)\sim2\log k\to\infty$.  Equivalently: $Q'$
lies in the complexified number-preserving algebra, the Levi component of
the parabolic-type stabilizer $P_\Omega$ of the vacuum line, on which the
Cartan norm is unbounded.

\subsection{No frame tames the map; the occupation is a state functional}

One might hope that a frequency-adapted symplectic frame $T_\w$ (the one
whose canonical vacuum is the physical two-field vacuum) turns the
divergent Cartan data into bounded physical squeezing.  It cannot:

\begin{lemma}[No frame tames the map]\label{lem:noframe}
For every symplectic frame $T$ (so that the Cartan data of
$T^{-1}S_+T$ are defined in the same group), in the Gram convention of
Section~\ref{sec:intro},
\begin{equation}
\big\|\mu\big(T^{-1}S_+T\big)\big\|_\infty
\ \ge\ \log\varrho\big((T^{-1}S_+T)^2\big)=\log\varrho(S_+^2)=r ,
\end{equation}
where $\varrho$ is the spectral radius: for any matrix $M$,
$\|\mu(M)\|_\infty=2\log s_{\max}(M)\ge2\log\varrho(M)
=\log\varrho(M^2)$, and similarity preserves the spectrum of the
positive $S_+$, which is $\{e^{\pm r/2}\}$.  Hence
$\|\mu(T_\w^{-1}S_+T_\w)\|\to\infty$ in the ultraviolet for every choice
of frames $T_\w$.  (Numerically the physical-frame Cartan data are in
fact close to $(r,r)$; the lemma is what the theorems below use.)
\end{lemma}

What is bounded is the \emph{state}: the selected vacuum, computed exactly
in \cite{companion2} as the Gaussian ground state of the PT mode
Hamiltonian, has occupation \eqref{eq:N} relative to the two-field vacuum,
saturating at $1/3$.  The two statements are compatible because particle
content depends on the transformation only through the vacuum ray it
produces --- a $G/P_\Omega$ datum --- while $\mu$ is a $G/K$ datum, and
the two are inequivalent quotients of $G$.

\begin{theorem}[Cartan--parabolic decoupling]\label{thm:decoupling}
Let $G=\Sp(2n,\C)$ act metaplectically on Fock space, and let
$q_K:G\to G/K$ and $q_\Omega:G\to G/P_\Omega$,
$q_\Omega(g)=[\hat g\,\Omega]$, be the two quotient maps.  Then:
(i) positive-metric geometry is a function of $q_K$ (the Cartan
projection); (ii) Gaussian vacuum rays are a function of $q_\Omega$;
(iii) $q_\Omega$ does not factor through $q_K$ --- there is no
$G$-equivariant $F$ with $q_\Omega=F\circ q_K$ --- since
$K\not\subseteq P_\Omega$ ($K\cap P_\Omega\cong U(n)\subsetneq
K\cong\Sp(n)$), so $q_\Omega$ is not constant on $K$-cosets; (iv) the
Cartan norm is unbounded on $P_\Omega$ (already
on its Levi component), so fibres of $q_\Omega$ contain elements of
arbitrarily large Cartan norm; consequently (v) particle content is not
a function of the Cartan projection, and no coercive lower bound
$\langle N\rangle\ge f(\|\mu\|)$ with $f(r)\to\infty$ holds in general.
(We do \emph{not} claim that no upper envelope
$\langle N\rangle\le f(\|\mu\|)$ exists; standard Bogoliubov estimates
bound particle number above by functions of the singular values, and on
bounded Cartan sets Gaussian occupation is continuous.  The decoupling
is one-sided: distortion does not \emph{force} excitation.)
The PU metric orbit realizes the extreme case:
$q_\Omega(S_+H)=\{[\Omega_{\mathrm{PT}}]\}$ is a single vacuum ray
(Theorem~\ref{thm:constancy}), while $q_K(S_+H)$ is unbounded, with
$\mu(S_+)=(r,r)$ divergent in the ultraviolet yet identically vanishing
canonical-frame excitation, and bounded physical excitation \eqref{eq:N}
despite Lemma~\ref{lem:noframe}.
\end{theorem}

The proof of (iv) is the identity \eqref{eq:beamsplitter} together with
$\mu(e^{(r/2)B})=(r,r)$; (v) follows because the exhibited family has
$\|\mu\|\to\infty$ with vanishing canonical-frame excitation, refuting
any coercive lower bound.  This corrects a natural but false expectation
--- used, e.g., to estimate Bogoliubov costs by $\sinh^2$ of Cartan
coordinates --- and it is the conceptual reason the theorems below must
be formulated at the level of vacuum functionals.

\begin{remark}[The geometric picture]\label{rem:wall}
In the Cartan-projection formulation of \cite{companion2}
(Section~\ref{sec:intro}), $\mu(S_+)=(r,r)$ lies on the $C_2$
Weyl-chamber wall for every $k$, and the ultraviolet limit
$r(k)\sim2\log k\to\infty$ is escape to infinity \emph{along the wall}
while remaining inside a single fibre of the vacuum-ray map $q_\Omega$.
This is the sharpest form of Cartan--parabolic decoupling: the metric
representative travels unboundedly in $G/K$ along a fixed direction
while the vacuum ray never moves.
\end{remark}

\section{Orbit constancy and the universal Fock obstruction}
\label{sec:obstruction}

\begin{lemma}[Eigenvector lemma]\label{lem:eigen}
Let $\ell_j$ be the physical annihilator covectors (the two-field vacuum
modes) and $X_j$ the stabilizer generators of \cite{companion2}.  Then
\begin{equation}
\ell_j^{\top}X_j=-i\,\ell_j^{\top},
\qquad\text{hence}\qquad
\ell_j^{\top}C(\theta_1,\theta_2)^{-1}=e^{i\theta_j}\,\ell_j^{\top}
\quad(\theta_j\in\C).
\end{equation}
\end{lemma}

\begin{theorem}[Metric-independence of the vacuum]\label{thm:constancy}
Call a positive metric \emph{admissible} if it belongs to the
classification $\eta'=\rho^{\dg}W\rho$ of \cite{companion1}, with
$W=f(N_1,N_2)>0$ densely defined with the polynomial CCR algebra of the
mode variables in its domain, and $\Omega$ an eigenvector of $W$
(automatic for $W=f(N_1,N_2)$: $W\Omega=w_0\Omega$, $w_0=f(0,0)>0$).
Then the selected physical vacuum ray, and the associated quasifree
functional on the polynomial CCR algebra, is the same for every
admissible positive metric:
(i) on the Gaussian orbit $\eta_C$,
$C\in SO(2,\C)^2$, the annihilation conditions defining the vacuum are
unchanged by Lemma~\ref{lem:eigen}; (ii) the metaplectic implementers of
the stabilizer are of the form
$\exp[\theta_1(N_1+\tfrac12)+\theta_2(N_2+\tfrac12)]$ and act on the
vacuum ray by scalars; (iii) for the full family, $W^{\pm1/2}$
also fixes the vacuum ray.  Hence $[\Omega_{\eta}]=[\Omega_{\eta_+}]$ for
every admissible $\eta$.
\end{theorem}

\begin{proof}[Ray constancy implies functional constancy]
Constancy of the ray alone does not suffice, because a metric change
alters both the physical inner product and the representatives of
observables; the functional is unchanged because the two changes cancel.
With $\rho'=W^{1/2}\rho$ and the functional defined by pullback on the
common polynomial CCR domain,
\begin{equation}
\omega_{\rho'}(A)
=\frac{\langle\Omega,\,W^{1/2}\rho\,A\,\rho^{-1}W^{-1/2}\,\Omega\rangle}
       {\langle\Omega,\Omega\rangle}
=w_0^{1/2}\,w_0^{-1/2}\,
 \frac{\langle\Omega,\,\rho A\rho^{-1}\,\Omega\rangle}
       {\langle\Omega,\Omega\rangle}
=\omega_{\rho}(A),
\end{equation}
using $W^{\pm1/2}\Omega=w_0^{\pm1/2}\Omega$ on both sides of the matrix
element.  The identity holds for every $A$ in the polynomial CCR
algebra; unbounded $W$ enters only through its action on $\Omega$ and on
the (analytic-vector) domain.
\end{proof}

Metric changes do alter the physical inner product and the
representatives of observables individually; the displayed cancellation
is what makes any quantity determined by the fixed vector relative to
the fixed two-field reference --- in particular the occupation
\eqref{eq:N} and the fidelity below --- metric-independent.

\begin{theorem}[Universal Fock obstruction]\label{thm:nogo}
Transport both constructions to the common auxiliary real CCR algebra of
the canonical variables $V=(x,y,p,q)$ per mode, with its standard
involution --- on this algebra both the two-field Fock vacuum and the
selected vacuum are positive Gaussian states (the selected vacuum is a
normalizable real Gaussian \cite{companion2}); the physical involution
differs by the PT rotation $z=iy$, which is genuinely complex, and this
is why positivity of the selected functional in \emph{physical} variables
is available only in the $\eta$-inner-product sense.  Then, for every
admissible positive metric, in a volume $V$ with cutoff
$\Lambda$: the relative quanta obey
$N(\Lambda)\sim C_d\,V\Lambda^d$ with $C_d>0$; the vacuum fidelity obeys
$|\langle\Omega_{2\text{-field}}|\Omega_{\mathrm{PU}}\rangle|^2
\le e^{-cV\Lambda^d}$ (per-mode fidelity $\to\sqrt3/2$
\cite{companion2}); and in the thermodynamic/continuum limit the
corresponding positive representations of the common auxiliary algebra
are disjoint, in every spatial dimension $d\ge1$.
Metric ambiguity cannot cure the inequivalence.
\end{theorem}

\begin{proof}
Combine Theorem~\ref{thm:constancy} with the exact modewise data
\eqref{eq:N} and the von Neumann infinite-tensor-product criterion, as in
\cite{companion2}; disjointness is meant for the stated common auxiliary
algebra, on which both states are positive under one involution.
\end{proof}

\section{The ultraviolet sector hierarchy}\label{sec:sectors}

Comparisons between selected vacua of different mass pairs are made in
common physical variables ($\gamma=1$ PU variables per mode; the
normalization map contributes only $O(k^{-2})$).  The exact modewise
Gaussian of \cite{companion2} gives, for pairs $(m_1,m_2)$ and
$(\bar m_1,\bar m_2)$, the fidelity expansion \eqref{eq:sector}; at
matched $\Sigma$ the expansion continues
\begin{equation}\label{eq:sector8}
1-f(k)\Big|_{\Sigma=\bar\Sigma}
=\frac{29}{576}\,\frac{(\delta^2-\bar\delta^2)^2}{k^{8}}+O(k^{-10}),
\qquad \delta^2-\bar\delta^2=\bar\Pi-\Pi ,
\end{equation}
i.e.\ at matched $\Sigma$ the $k^{-6}$ coefficient also vanishes and the
first remaining obstruction occurs at order $k^{-8}$.

\begin{theorem}[Ultraviolet sector hierarchy]\label{thm:sectors}
Let $d$ be the spatial dimension, and hold the infrared behaviour of the
compared vacua in a common regular class (in particular, strictly
positive masses).  The relative ultraviolet Bogoliubov transformation is
Hilbert--Schmidt --- diagnosed by
$\int^\Lambda k^{d-1}(1-f(k))\,dk<\infty$, which for nearby ultraviolet
Gaussians is proportional to the squared Bogoliubov coefficient ---
precisely according to:
\begin{center}
\begin{tabular}{c|c}
spatial dimension & UV sector invariant of selected vacua\\
\hline
$d\le3$ & none (no UV invariant)\\
$4\le d<8$ & $\Sigma=m_1^2+m_2^2$\\
$d\ge8$ & $(\Sigma,\Pi)$, i.e.\ the unordered mass pair
\end{tabular}
\end{center}
The hierarchy terminates: the two symmetric invariants determine the
pair up to interchange $(m_1^2,m_2^2)\sim(m_2^2,m_1^2)$.
At the critical dimensions the obstruction is logarithmic rather than
power-law: at $d=4$ (respectively $d=8$) the divergent sum for
$\Sigma\neq\bar\Sigma$ (respectively $\Pi\neq\bar\Pi$ at matched
$\Sigma$) grows as $\log\Lambda$, so the critical dimensions belong to
the labeled regimes as indicated.  This theorem classifies ultraviolet
(local) sectors; global equivalence additionally requires an infrared
analysis, as follows.
\end{theorem}

\begin{remark}[The $\Box^2$ anchor is an infrared question, left open]
\label{rem:anchor}
The vacuum-overlap criterion is an ultraviolet diagnostic only: $1-f$ is
bounded by $1$, so its infrared integral always converges, whereas the
global Shale criterion involves the relative occupation
$N_{\mathrm{rel}}$, which need not stay bounded.  Against the doubly
massless selected vacuum ($m_1=m_2=0$; modewise Gaussian
$A_0=(2k,k^2;k^2,2k^3)$), a massive selected vacuum has (machine-verified)
\begin{equation}
N_{\mathrm{rel}}(k)\ \sim\ \frac{(m_1+m_2)\,m_1m_2}{6\,k^{3}}
\qquad(k\to0),
\end{equation}
so $\int_0k^{d-1}N_{\mathrm{rel}}(k)\,dk$ diverges precisely for
$d\le3$: the dimensions with the most benign ultraviolet behaviour are
exactly those in which the massless anchor carries the strongest
potential infrared obstruction.  Whether the $\Box^2$ vacuum lies in the
same \emph{global} Fock/GNS sector as the massive selected vacua
therefore requires a separate infrared (Shale/Araki--Yamagami) analysis,
which we do not perform here.  For strictly positive masses the infrared
is uniformly regular and the hierarchy of Theorem~\ref{thm:sectors}
stands as stated.
\end{remark}

At the resonant boundary four kinds of regularity must be
distinguished, and they are logically independent.
\emph{Distributional regularity}: the selected two-point functional is
analytic in $\Delta$ across $\Delta=0$ (modewise, $1-f\sim\epsilon^4$
for $m_{1,2}=m\pm\epsilon$; the Wightman functions of
Section~\ref{sec:bridge} converge smoothly), and the occupation limits
commute ($\langle N\rangle\equiv\tfrac13$ at $\Delta=0$ for every $k$).
\emph{Representation regularity at fixed positive mass}: the confluence
$m_1,m_2\to m>0$ introduces no new infrared singularity, and for
$d\le3$ the massive resonant vacuum lies in the same ultraviolet-dressed
sector as every nondegenerate selected vacuum with the same infrared
class (Theorem~\ref{thm:sectors}).  \emph{The doubly massless corner}
$m=0$ is a separate boundary governed by Remark~\ref{rem:anchor} and is
not settled by the ultraviolet hierarchy.
\emph{Dynamical and metric singularity}: as $\Delta\to0$ the metric
operator degenerates ($\alpha,\beta\to\infty$), the similarity
transformation diverges ($S_+\sim\Delta^{-1/2}$), and the dynamics becomes
nondiagonalizable (Jordan blocks; no positive invariant form
\cite{companion1}).  In summary,
\begin{equation}
\boxed{\ \text{resonance is not a vacuum singularity; it is a metric and
dynamical singularity.}\ }
\end{equation}

\section{The spectral bridge to the Krein quantization}\label{sec:bridge}

The physical Wightman function of the selected vacuum follows from the
exactly solvable mode: with
$\widetilde V=S_+V$ evolving under the diagonalized Hamiltonian,
$W(t)=S_+e^{tA_0}\Gamma S_+^{\top}$ pulled back to the field variable
$z=iy$ gives, exactly,
\begin{equation}\label{eq:Wmode}
W_{zz}(t)
=\frac{1}{\Delta}
\left[\frac{e^{-i\w_1t}}{2\w_1}-\frac{e^{-i\w_2t}}{2\w_2}\right],
\qquad
W_{zz}(t)\Big|_{\Delta\to0}
=-\,\frac{(1+i\w t)\,e^{-i\w t}}{4\w^3}.
\end{equation}
Both expressions contain only positive frequencies; the first is the
momentum-space spectral function \eqref{eq:Wspectral}, the second its
confluent ($\delta'$) limit.  The commutator parts reproduce the classical
fourth-order Green-function difference
$E_P=[\sin\w_2t/\w_2-\sin\w_1t/\w_1]/\Delta$ with the normalization
$E_P(0)=E_P'(0)=E_P''(0)=0$, $E_P'''(0)=1$, and confluent limit
$(\sin\w t-\w t\cos\w t)/(2\w^3)$.

\begin{theorem}[Spectral bridge]\label{thm:bridge}
Fix the action sign and Fourier convention of \eqref{eq:action}
($\varepsilon=+1$, $E_P'''(0)=+1$).  Three statements, in decreasing
generality:
\begin{enumerate}
\item \emph{Split masses.}  For $m_1\ne m_2$, the selected PU covariance
is the positive-energy spectral splitting of the fourth-order
commutator, \eqref{eq:Wspectral}: metric selection and the spectral
condition determine the same complex quasifree covariance.
\item \emph{General resonance.}  For $m_1,m_2\to m$ (any $m\ge0$), the
confluent limit is
\begin{equation}
\widetilde W^+_{mm}(p)=-\varepsilon\,\theta(p^0)\,\delta'(p^2-m^2)
=\varepsilon\,\partial_{m^2}\big[\theta(p^0)\delta(p^2-m^2)\big],
\end{equation}
a \emph{massive Bateman--Turok-type} Krein/Jordan covariance, realized
by a Jordan mode doublet with null pairing
$[a_1,a_2^{\dg}]=[a_2,a_1^{\dg}]\ne0$.  This massive family is not the
theory treated in \cite{BT2026}.
\item \emph{Doubly massless point.}  At $m=0$, and after matching the
infrared extension, the confluent covariance agrees with the
Bateman--Turok covariance
$\widetilde W_{\mathrm{BT}}(p)=\theta(p^0)\delta_1'(p^2)$,
$\delta_1'(p^2)=-\partial_{m^2}\delta(p^2-m^2)|_{m^2=0}$
\cite{BT2026}, up to the overall action-sign reversal
$\varepsilon\to-\varepsilon$: in our $\varepsilon=+1$ convention the two
covariances are negatives of each other, and they coincide in the
perfect-square convention $\varepsilon=-1$.
\end{enumerate}
For $\Delta>0$ the covariance is realized through a positive
pseudo-Hermitian completion ($\eta$-inner product); at $\Delta=0$ it is
realized through the resonant Krein/Jordan completion with ghost parity.
Equality throughout concerns the complex quasifree covariance
(Conventions); we do not assert positivity with respect to a common
involution.
\end{theorem}

\begin{proof}
Statement 1 is the computation \eqref{eq:Wmode} (machine-verified,
including the commutator normalization $E_P'''(0)=+\varepsilon$, which is
what ties the covariance to the action sign).  Statement 2 is the
machine-verified confluent limit of \eqref{eq:Wmode}: the divided
difference converges to $+\partial_{m^2}W^+_m
=-(1+i\w t)e^{-i\w t}/(4\w^3)$ per mode.  Statement 3 is then the
identity $\delta_1'(p^2)=-\partial_{m^2}\delta(p^2-m^2)|_0$: the
Bateman--Turok mode function is the negative of ours, which is exactly
the sign of their perfect-square action.  Quasifree functionals are
determined by their two-point distributions, so these mode-level
identities identify the covariances as stated.
\end{proof}

\begin{remark}
(i) The theorem upgrades the resonant-boundary picture: the Krein
formulation is not the residue of a positive theory destroyed by an
ultraviolet particle catastrophe (Section~\ref{sec:decoupling} shows no
such catastrophe exists); it is a different completion of the
\emph{same}, analytically continued covariance, adapted to Jordan
dynamics.  (ii) At $\Delta>0$ the completion ambiguity is invisible in
expectation values but material for which operators are Hermitian.  At
$\Delta=0$ no nondegenerate positive-definite invariant metric exists
(the Jordan obstruction of \cite{companion1}), and the Krein
construction is the canonical nondegenerate spectral completion
considered here; singular or intrinsically Jordan $PT$-symmetric
boundary realizations in the sense of Bender and Mannheim
\cite{BM2008PRD} are distinct constructions not classified by this
theorem.  (iii) What
remains genuinely open at the resonant point is the treatment of the
generalized (Jordan) excitations, where the two completions differ by
construction; ghost parity governs precisely that sector \cite{BT2026}.
\end{remark}

\section{Hadamard structure of the selected vacuum}\label{sec:hadamard}

Let $W^+_m$ denote the Minkowski Klein--Gordon positive-frequency
two-point function and define the divided difference and its confluent
limit
\begin{equation}
W^+_{12}=\frac{W^+_{m_1}-W^+_{m_2}}{m_1^2-m_2^2},
\qquad
W^+_{mm}=+\,\partial_{m^2}W^+_m
\end{equation}
(the confluent limit of a divided difference is the $+$derivative;
machine-verified against \eqref{eq:Wmode}), which by
Theorem~\ref{thm:bridge} are the selected two-point functions in the
$\varepsilon=+1$ convention.

\begin{theorem}[Fourth-order Hadamard property]\label{thm:hadamard}
Let $P=(\Box+m_1^2)(\Box+m_2^2)$, $d=3$ spatial dimensions.  Then:
\begin{enumerate}
\item $P_xW^+_{12}=P_yW^+_{12}=0$, and
$W^+_{12}(x,y)-W^+_{12}(y,x)=iE_P(x,y)$ with
$E_P=(E_{m_1}-E_{m_2})/\Delta$ the fourth-order commutator function;
\item $\WF(W^+_{12})=\mathcal C^+$, the standard positive-frequency
Hadamard wavefront set; the multiplicity of the characteristic polynomial
changes the order of the light-cone singularity, not its wavefront
directions;
\item the short-distance singularity is logarithmic, in the Hadamard
form
\begin{equation}
W^+_{12}(x,y)=V_{12}(x,y)\,\log\sigma_+(x,y)+H_{12}(x,y),
\end{equation}
with $V_{12}$ and $H_{12}$ smooth and universal (mass-independent)
nonzero coincidence value $V_{12}(x,x)=\frac{1}{16\pi^2}$ (in the
$\log\rho^2$ normalization; $\frac1{8\pi^2}$ for $\log\rho$): with
$W_m^+(x)=\frac{1}{4\pi^2}\frac{mK_1(m\rho)}{\rho}$,
$\rho^2=-x^2+i0x^0$, the mass-independent $\frac{1}{4\pi^2\rho^2}$ term
cancels in the divided difference and, in flat coordinates,
\begin{equation}
W^+_{12}(x)=\frac{1}{8\pi^2}\log\rho
+\frac{\Sigma}{64\pi^2}\,\rho^2\log\rho+\dots+C^\infty ,
\end{equation}
where the displayed subleading term (coefficient machine-verified) is
milder than the leading logarithm but \emph{not} smooth for nonzero
masses: the universal statement concerns the coincidence value of the
smooth logarithmic coefficient, not constancy of the coefficient modulo
smooth terms.  The confluent case has the identical leading coefficient;
\item in the doubly massless case $P=\Box^2$ the infrared extension of
$\theta(p^0)\delta'(p^2)$ requires a scale; different admissible scales
change $W$ by a smooth constant, affecting neither the wavefront set nor
the local Hadamard property --- the bridge of Theorem~\ref{thm:bridge}
presumes matched extension conventions;
\item any two fourth-order spectral Hadamard functionals with the same
commutator differ by a smooth bisolution; the fourth-order parametrix
$H_P=(H_{m_1}-H_{m_2})/\Delta$ (confluently $\partial_{m^2}H_m$), whose
full smooth coefficient multiplies $\log\sigma_+$ as in (3), has
$W^+_{12}-H_P$ smooth, and provides the local microlocal input for
developing Wick powers $:\!\phi^n\!:_{H_P}$ and perturbative
time-ordered products.  (The Hollands--Wald and Brunetti--Fredenhagen
constructions are formulated for hyperbolic operators and positive
Hadamard states; their adaptation to a fourth-order operator with
repeated characteristics at resonance, an indefinite quasifree
functional, and a nonstandard involution is a genuine further step, not
claimed here.)
\end{enumerate}
\end{theorem}

\begin{proof}[Proof sketch]
(1) is verified exactly modewise (both branches and the confluent Jordan
mode), with the Green normalization $E_P'''(0)=1$.  For (2), Fourier
support in the closed forward cone alone does not give the
position-space wavefront relation; the inclusion
$\WF\subseteq\mathcal C^+$ follows because divided differences and
$\partial_{m^2}$ of the smooth family $m^2\mapsto W^+_m$ preserve the
Hadamard wavefront bound of the Klein--Gordon factors.  For the reverse
inclusion, use the boundary-value form of $\log\sigma_+$ (equivalently
the explicit Hadamard expansion of (3)): the mass-independent
$1/\sigma$ singularities cancel, which \emph{lowers the singular order}
but removes no covector directions, because the smooth coefficient
$V_{12}$ has nonzero diagonal value $1/(16\pi^2)$; hence no component of
$\mathcal C^+$ disappears, and the result is a \emph{fourth-order
microlocal spectrum condition} (we avoid ``Hadamard state'', which
normally presupposes positivity relative to a fixed involution).  (3) is
the Bessel expansion
$mK_1(m\rho)/\rho=\rho^{-2}+\tfrac{m^2}{2}\log(m\rho/2)+\dots$: the
$\rho^{-2}$ terms cancel; the $\log\rho$ coefficient of the divided
difference is $\big[\tfrac{m_1^2}{2}-\tfrac{m_2^2}{2}\big]/\Delta\cdot
\tfrac{1}{4\pi^2}=\tfrac{1}{8\pi^2}$ and the $\rho^2\log\rho$
coefficient is $\Sigma/(64\pi^2)$ (both machine-verified, as is the
confluent case).  (4) is the observation that the scale enters only
through $\log(\mu\rho)$, so a change of scale adds
$\log(\mu/\bar\mu)/(8\pi^2)$, a constant.  (5) follows from (1)--(3) by
the usual difference argument: the difference of two such functionals is
a bisolution with vanishing antisymmetric part and empty wavefront set.
\end{proof}

\begin{proposition}[$(\pm)$Hadamard split structure]\label{prop:split}
Let $\phi_1=(\Box+m_2^2)\phi/(m_2^2-m_1^2)$ and
$\phi_2=(\Box+m_1^2)\phi/(m_1^2-m_2^2)$ be the local partial-fraction
fields.  In the selected vacuum, exactly,
\begin{equation}
W_{\phi_1\phi_1}=+\frac{W^+_{m_1}}{\Delta},
\qquad
W_{\phi_2\phi_2}=-\frac{W^+_{m_2}}{\Delta},
\qquad
W_{\phi_1\phi_2}=0 .
\end{equation}
\end{proposition}

The ghost branch carries a Klein--Gordon--Hadamard \emph{distribution
with negative Krein signature} --- not a positive KG state under the
ordinary KG $*$-structure: the Krein signature is intrinsic to the
selected functional, while its local singularity structure is standard on
both branches.  This makes the summary statement precise:

\begin{quote}
\emph{The universal selected vacuum is a fourth-order Hadamard quasifree
functional.  Its apparent infinite particle content relative to the split
two-field vacuum reflects inequivalent complex structures --- the
positive-positive versus positive-negative pairing of the two branches
--- not pathological local ultraviolet singularities.}
\end{quote}

The logarithmic singularity is the expected one for a dimension-zero
field in four spacetime dimensions and is milder than the Klein--Gordon
$1/\sigma$; local Wick composites relative to $H_P$ are well-defined at
the level of the subtraction, and derivative observables (in particular
the stress tensor) require the usual local subtractions relative to
$H_P$, with no evident ultraviolet obstruction.  A systematic treatment
along the lines of \cite{Radzikowski,BrunettiFredenhagen,HollandsWald}
for the fourth-order operator --- which would have to accommodate
repeated characteristics at resonance, indefiniteness, and the
nonstandard involution --- now has its local microlocal input in place.

\section{Discussion}\label{sec:discussion}

\paragraph{What the three geometries settle.}
The separation \eqref{eq:threegeom} resolves several tensions in the
literature at once.  The divergent rapidity $r(k)\sim2\log k$ is real,
but it is metric geometry: it makes the phase-space (observable) metric a
Sobolev pair $H^1\oplus H^{-1}$ \cite{companion2} and the Dyson operator
unboundedly non-unitary in every frame --- yet it creates no quanta,
because it acts inside the vacuum-line stabilizer.  The physical
inequivalence to the naive two-field Fock space is vacuum geometry: an
$O(1)$-per-mode, metric-independent statement
(Theorems~\ref{thm:constancy},~\ref{thm:nogo}).  The resonant
$\Delta\to0$ singularity is dynamical geometry: Jordan blocks and the
death of the positive metric, with the vacuum analytic throughout
(Theorem~\ref{thm:sectors} ff.).  Conflating any two of these produces
false conclusions --- e.g.\ ultraviolet particle catastrophes
($\Lambda^{d+2}$ scalings) or mass-splitting superselection sectors, both
of which are refuted by the exact computations here.

\paragraph{Relation to Bender--Mannheim and Bateman--Turok.}
Theorem~\ref{thm:bridge} shows the two programs select the same complex
covariance,
from opposite directions: positivity of the inner product modewise
(Bender--Mannheim), spectral positivity of the energy covariantly
(Bateman--Turok).  The difference is the completion, and the difference
matters exactly where the completions are forced apart: at the resonant
point, where positivity fails (the Jordan no-go of \cite{companion1}) and
ghost parity takes over the excitation sector.  A sharp remaining
question is the representation-theoretic status of the generalized Jordan
excitations at $\Delta=0$: the vacuum sector is common, the one-particle
completions are not, and the interacting perfect-square theory of
\cite{BT2026} lives precisely there.

The division of labor should be stated plainly.  The Krein quantization
itself, the hidden ghost-parity symmetry, the two-field $O(1,1)$
embedding, and the positivity of transition probabilities on an
indefinite state space are constructions of Bateman and Turok
\cite{BT2026}; none of them is claimed here.  What this paper
contributes is the identification of the common \emph{free complex
covariance} underlying both programs and the precise determination of
which additional data --- involution, observable algebra, and completion
--- distinguish them.  Two consequent cautions.  First, a common free
complex covariance does not imply an identical interacting quantum
theory: the companion interaction paper \cite{companion5} shows the two
completions \emph{separate} under interaction --- the canonical positive
metric is obstructed by on-shell branch-changing conversion, while the
Krein/ghost-parity structure survives the same interaction exactly.
Second, the ``Krein boundary'' of the title is a statement about the
free covariance: the confluent covariance naturally admits an
indefinite/Krein completion; it does not by itself determine the
interacting Bateman--Turok theory, whose $O(1,1)$, null-state, and
ghost-symmetry structure is additional data beyond the covariance.

\paragraph{Toward quadratic gravity.}
The scalar results transfer directly to genuine PU blocks: the two
transverse-traceless helicity blocks of scalar-free Einstein--Weyl
gravity, and scalaron blocks when present.  They do \emph{not} transfer
independently to every spin-two polarization: after diffeomorphism
reduction the physical phase space stratifies as
$\Gamma_{\mathrm{phys}}\cong2\Gamma_{\mathrm{PU}}\oplus2\Gamma_-
\oplus\Gamma_-$ --- only helicity $\pm2$ survives as genuine PU pairs,
the massive helicities $\pm1,0$ are unpaired second-order ghosts
(their massless partners lie in the diffeomorphism orbit), and Lorentz
covariance then forces all five massive polarizations into one joint
real-form choice \cite{companion4}.  Where the transfer applies, the
physical vacuum is spectral, satisfies the fourth-order microlocal
spectrum condition, is metric-independent, and is not normal to the
naive two-field (graviton plus ghost) Fock representation; PT-based
unitarity claims should accordingly be phrased in the dressed
representation, and interacting constructions must fix the
representation before the perturbative expansion.  We leave the
interacting analysis, the stress-tensor subtractions relative to $H_P$,
and curved backgrounds to future work.

\begin{thebibliography}{12}

\bibitem{BM2008PRL}
C.~M.~Bender and P.~D.~Mannheim,
\emph{No-ghost theorem for the fourth-order derivative Pais--Uhlenbeck
oscillator model},
Phys.\ Rev.\ Lett.\ \textbf{100} (2008) 110402.

\bibitem{BM2008PRD}
C.~M.~Bender and P.~D.~Mannheim,
\emph{Exactly solvable PT-symmetric Hamiltonian having no Hermitian
counterpart},
Phys.\ Rev.\ D \textbf{78} (2008) 025022, arXiv:0804.4190.

\bibitem{BT2026}
S.~Bateman and N.~Turok,
\emph{Escape from Ostrogradsky via hidden ghost parity},
arXiv:2607.00096.

\bibitem{companion1}
GPT-5.6.sol and A.~A.~Palm,
\emph{Canonical Positive Symplectic Diagonalization of the
Pais--Uhlenbeck Oscillator},
Paper~I of the pure-Weyl gravity programme, 2026;
file \path{paper/01-symplectic-diagonalization.pdf} in the repository
\url{https://github.com/area9innovation/weyl-gravity}.

\bibitem{companion2}
GPT-5.6.sol and A.~A.~Palm,
\emph{The Pais--Uhlenbeck Metric as a Minimum-Distortion Principle,
and the Representation Problem for the Fourth-Order Field},
Paper~II of the pure-Weyl gravity programme, 2026;
file \path{paper/02-variational-fock.pdf} in the repository
\url{https://github.com/area9innovation/weyl-gravity}.

\bibitem{companion4}
GPT-5.6.sol and A.~A.~Palm,
\emph{Gauge Reduction and the Completion Problem in Fourth-Order Gravity:
PU Pairing, Covariant Real Forms, and the Conformal Jordan Boundary},
Paper~IV of the pure-Weyl gravity programme, 2026;
file \path{paper/04-fourth-order-gravity.pdf} in the repository
\url{https://github.com/area9innovation/weyl-gravity}.

\bibitem{companion5}
GPT-5.6.sol and A.~A.~Palm,
\emph{Interaction Obstructions, Resonant PT Breaking,
and Doubled Jordan Symmetry in Fourth-Order Theories},
Paper~V of the pure-Weyl gravity programme, 2026;
file \path{paper/05-interaction-obstructions.pdf} in the repository
\url{https://github.com/area9innovation/weyl-gravity}.

\bibitem{Mostafazadeh2010}
A.~Mostafazadeh,
\emph{Pseudo-Hermitian representation of quantum mechanics},
Int.\ J.\ Geom.\ Meth.\ Mod.\ Phys.\ \textbf{7} (2010) 1191.

\bibitem{vN1939}
J.~von Neumann,
\emph{On infinite direct products},
Compos.\ Math.\ \textbf{6} (1939) 1.

\bibitem{Radzikowski}
M.~J.~Radzikowski,
\emph{Micro-local approach to the Hadamard condition in quantum field
theory on curved space-time},
Commun.\ Math.\ Phys.\ \textbf{179} (1996) 529,
\href{https://doi.org/10.1007/bf02100096}{doi:10.1007/bf02100096}.

\bibitem{BrunettiFredenhagen}
R.~Brunetti and K.~Fredenhagen,
\emph{Microlocal analysis and interacting quantum field theories},
Commun.\ Math.\ Phys.\ \textbf{208} (2000) 623.

\bibitem{HollandsWald}
S.~Hollands and R.~M.~Wald,
\emph{Local Wick polynomials and time ordered products of quantum fields
in curved spacetime},
Commun.\ Math.\ Phys.\ \textbf{223} (2001) 289,
\href{https://doi.org/10.1007/s002200100540}{doi:10.1007/s002200100540}.

\bibitem{Holdom}
B.~Holdom,
\emph{UV-complete 4-derivative scalar field theory},
Nucl.\ Phys.\ B \textbf{1000} (2024) 116472,
\href{https://doi.org/10.1016/j.nuclphysb.2024.116472}
{doi:10.1016/j.nuclphysb.2024.116472};
\emph{Making sense of ghosts},
Nucl.\ Phys.\ B \textbf{1008} (2024) 116696,
\href{https://doi.org/10.1016/j.nuclphysb.2024.116696}
{doi:10.1016/j.nuclphysb.2024.116696}.

\bibitem{ABHT}
M.~Anderson, S.~Bateman, F.~Herzog and N.~Turok,
\emph{Renormalization of a four-derivative theory} and
\emph{Conformally flat limit of quadratic gravity}, 2026.
Both are announced as forthcoming and were unpublished as of July 2026; they
are cited as ``to appear'' in refs.~[25] and~[27] of \cite{BT2026}, which is
the only public record of them we have located.

\end{thebibliography}

\end{document}
